\section*{Summary}

\pending{\textbf{Preview (4 November 2026).} Six runs are still queued on
merlin7. The one that matters is \ffn{} FE3 $r=2$ on IAEA-3D, which
completes the IAEA reference; until it is in, the IAEA results rest on a
provisional reference (Section~\ref{sec:steady}). The other five add points
to the figures and the run tables and cannot change a time-to-accuracy
entry. Text in this colour depends on the pending runs. Every table and
figure is regenerated from the data.}

This report times \garnet{} against \ffn{} on the same multigroup diffusion
problems, solved to the same accuracy:
\begin{itemize}
\item IAEA-3D;
\item the 3-D Langenbuch core at $t=0$;
\item the 11-group hex-homogenized HP-MR;
\item a $+0.5$\,\$ reactivity insertion in the HP-MR.
\end{itemize}
Both codes run on one core of the same AMD EPYC 7742 node, and \garnet{}
also on one GH200 GPU. Every run is scored against a common reference. For
each accuracy target, the comparison takes each platform's cheapest
measured run that meets it (Table~\ref{tab:tta}).

\paragraph{Findings.}
\begin{itemize}
\item \textbf{Steady state, one core: \ffn{} is faster.} At 10~pcm in $k$
  and 1\% in every node power, its degree-2 and degree-3 elements need 8 to
  20 times fewer unknowns than \garnet's second-order finite volumes.
  \begin{itemize}
  \item HP-MR: 89~s for \ffn{} against 2000~s for \garnet.
  \item Langenbuch: 45~s against 66~s.
  \item 1~pcm and 0.1\%: only \ffn{} reaches it (Langenbuch 434~s,
    HP-MR 124~s).
  \end{itemize}
\item \textbf{Steady state, GH200: \garnet{} is fastest wherever its meshes
  reach the target.} HP-MR in 7.0~s, 13 times faster than \ffn{} on one
  core; Langenbuch in 0.35~s, 130 times faster.
\item \textbf{IAEA-3D: no \garnet{} run meets the 1\% node-power target.}
  At $r=16$ (37.5 million unknowns) its $k$ is within 10~pcm. Its worst
  node, at the corner of the fuel next to the radial and axial reflectors,
  is still \pending{2.6\%} off. \ffn{} meets the target in \pending{658~s}.
\item \textbf{Transient: \garnet{} is faster on the same core.} To keep the
  power error below $10^{-3}$:
  \begin{itemize}
  \item \ffn{}: 1262~s on one core;
  \item \garnet{}: 127~s on one core (10 times faster) and 17~s on the
    GH200 (76 times).
  \end{itemize}
  At $10^{-2}$ the two codes are even on one core (107 and 100~s).
\item \textbf{Why the transient differs: time integration.} \garnet's BDF2
  is second order; \ffn's two usable schemes are first order. \ffn's default
  scheme, Implicit-Exponential, does not hold the unperturbed state on this
  input.
\end{itemize}

\paragraph{Scope.}
\begin{itemize}
\item \ffn{} runs on one core only: its MPI path does not work yet
  (Section~\ref{sec:codes}). The GH200 figures are therefore against one
  core, not node against node.
\item The transient reference is \garnet's extrapolated GH200 solution.
  \ffn's own runs converge onto it independently.
\item The HP-MR drums are withdrawn. A smeared B$_4$C drum is black to
  thermal neutrons, and neither code resolves it.
\item Posing the problems identically in both codes uncovered five \ffn{}
  issues, listed in Section~\ref{sec:discussion}.
\end{itemize}
